\\maketitle
\\begin{abstract}
    本文针对2024年CUMCM B题碳排放预测与减排策略问题，建立了基于Transformer的碳排放预测与多目标减排优化模型。首先对历史碳排放数据进行特征工程，提取GDP、能源结构、产业结构等多维度影响因素；其次构建Transformer预测模型，实现未来10年碳排放趋势预测；最后建立多目标优化模型，以最小化减排成本和最大化经济发展为目标，采用MOEA/D算法求解。
    
    \\textbf{针对问题一}，我们建立了碳排放影响因素分析模型。通过Pearson相关系数分析，识别出GDP增长率、第二产业占比、煤炭消费占比为主要影响因素，解释了85\\%的碳排放方差。
    
    \\textbf{针对问题二}，我们构建了Transformer预测模型。采用多头自注意力机制捕捉时间序列长距离依赖关系，位置编码编码时序信息：
    $$Attention(Q,K,V) = softmax(\\frac{QK^T}{\\sqrt{d_k}})V$$
    预测结果显示，未来10年碳排放将呈现先增后降趋势，2030年前后达到峰值。
    
    \\textbf{针对问题三}，我们建立了多目标减排优化模型。考虑减排成本、经济发展、能源安全三个相互冲突的目标，采用MOEA/D算法求解Pareto最优解集，并通过熵权-TOPSIS方法进行方案排序。
    
    本研究为政府制定碳达峰碳中和政策提供了科学依据和决策支持。
    
    \\keywords{碳排放预测\quad Transformer\quad 多目标优化\quad MOEA/D\quad 碳达峰}
\\end{abstract}
